\subsection{Лекция 5 (17.10.26) МКО для радиально симметричной задачи}
В программе из файла \ename{poisson_log_fvm_test.cpp} рассматривается радиально симметричная задача в постановке
\begin{align*}
&\frac1r\dfr{}{r}\left(\lambda(r) r\dfr{u}{r}\right) = 0, \\
&u(r_0) = a, \quad u(r_1) = b,\\
&\lambda(r) = \begin{cases}
\lambda_0, & r < r_k,\\
\lambda_1, & r \geq r_k.\\
\end{cases}
\end{align*}
Общий вид решения этого уравнения примет вид
\begin{equation*}
u(r) = 
\begin{cases}
u_{0} = A_0 \ln(r) + B_0, & r < r_k\\
u_{1} = A_1 \ln(r) + B_1, & r \geq r_k.
\end{cases}
\end{equation*}
Четыре неизвестные константы $A_{0,1}, B_{0,1}$ находятся из двух граничных условий плюс двух условий
(равенство значений и потоков) на границе скачка коэффициента диффузии:
\begin{equation}
\label{eq:fvm_radial_partial_lambda_conditions}
\begin{array}{l}
r = r_0: \quad u_0 = a;\\
r = r_1: \quad u_1 = b;\\
r = r_k: \quad u_0 = u_1, \quad \lambda_0\dsfr{u_0}{r} = \lambda_1\dsfr{u_1}{r}.
\end{array}
\end{equation}
Использовались параметры
\begin{equation*}
r_0 = 0.05, \quad r_1 = r_0 + 1.0, \quad r_k = r_0 + 0.5, \quad a = 1, \quad b = 0, \quad \lambda_0 = 10, \quad \lambda_1 = 1.
\end{equation*}

В тесте \cvar{poisson-log-fvm-2d}
эта задача решена в двумерных декартовых координатах методом конечных объёмов на неструктурированной сетке.
Отталкиваясь от данной реализации необходимо:

\begin{enumerate}
\item Посчитать нормы на сгущающихся сетках и нарисовать графики сходимости.
\item Изменить вариант расчёта $\lambda$ на грани со среднеарифметического на среднегармонический. Построить график сеточной сходимости.
\item Добавить учёт логарифмического профиля при вычислении производной на грани $r=r_0$. Построить график сеточной сходимости.
\item Решить ту же самую задачу в одномерной радиально симметричной постановке.
      Рассмотреть варианты:
      \begin{itemize}
      \item среднеарифметическое $\lambda$ без поправки на логарифмический профиль,
      \item среднеарифметическое $\lambda$ с поправкой на логарифмический профиль,
      \item среднегармоническое $\lambda$ без поправкой на логарифмический профиль,
      \item среднегармоническое $\lambda$ с поправкой на логарифмический профиль.
      \end{itemize}
      Для каждого случая построить график сеточной сходимости
\item Проиллюстрировать полученные двумерные и одномерные решения с использованием сходных по разбиению
      (достаточно грубых) сеток.
      Сравнить численные решения с точным на едином одномерном графике.
\end{enumerate}

\paragraph{Рекомендации}
\begin{itemize}
\item
Для построение неструктурированных сеток в кольце $[r_0, r_1]$ для расчёта сеточной сходимости разной размерности
в пп.1--3 следует использовать скрипт \ename{test_data/radial.py} согласно \secref{sec:prog_grid}.

\item
Изменения для вычисления $\lambda$ на гранях (п.2) следует вносить в метод \cvar{build_lambda_face} (цикл по внутренним граням).
Использовать формулу \cref{eq:fvm_lambda_ij_2}.

\item
Учёт логарифмического профиля (п.3) также следует проводить через поправку на значение $\lambda$
на граничной грани при $r=r_0$ в методе \cvar{build_lambda_face} (цикл по граничным граням, явно выделяя грань с центром около $r=r_0$) согласно формуле \cref{eq:fvm_ln_singularity_hij}.
Для расчёта этой формулы понадобятся значения радиусов в центре приграничной ячейке ($r_i$) и  в центре граничной грани ($r_j$),
а также расстояние между этими центрами $h_{ij}$.
Эти значения можно расчитать получив соответсвующие центры как
\begin{cppcode}
Point pi = grid_->cell_center(icell);  // центр приграничной ячейки
Point pj = grid_->face_center(iface);  // центр грани
\end{cppcode}

\item 
Как отмечено в \secref{sec:fvm_radsym}, для метода МКО решение радиально симметричной
задачи в радиальных координатах отличается от решения в декартовых координатах
только способом вычисления геометрических характеристик сетки.
Поэтому для выполнения п.4 будет удобно ввести новый класс одномерной сетки \cvar{RadialGrid1D},
который отнаследовать от имеющегося класса \cvar{Grid1D} и переопределить
методы вычисления объёма ячейки \cref{eq:fvm_radsym_E_1d}, площади грани \cref{eq:fvm_radsym_gamma_1d} и центра ячейки \cref{eq:fvm_radsym_Ec_1d}:
\begin{cppcode}
class RadialGrid1D: public Grid1D{
public:
    RadialGrid1D(double r0, double r1, size_t n_cells): Grid1D(r0, r1, n_cells){}
    double face_area(size_t iface) const override {
        // TODO
    }
    double cell_volume(size_t icell) const override {
        // TODO
    }
    Point cell_center(size_t icell) const override {
        // TODO: return Point{rc, 0, 0}
    }
};
\end{cppcode}
Далее создавая эту сетку и подставляя её в имеющийся решатель
\begin{cppcode}
auto grid = std::make_shared<RadialGrid1D>(0.05, 1.05, n_cells);
Problem worker(grid, test_data);
\end{cppcode}
можно получать одномерное численное решение и норму его отклонения от истинного.
При этом алгоритмы расчёта $\lambda$ на гранях (написанные для пп.2--3) останутся неизменными.
Для этого случая следует подумать над оптимизацией количества внутренних итераций для поправки на скошенность
(аргумент \cvar{n_iters} в методе \cvar{solve}).

\item
Для отрисовки всех решений на едином одномерном графике (п.5) нужно открыть в паравью
все необходимые для отрисовки поля (точное, одномерные и двумерные), 
и далее следуя указаниям из \secref{sec:paraview-plot-over-line}
в каждом создать линию вдоль которой нужно построить решение (для 1D она будет совпадать c областью решения,
в 2D -- любая радиальная линия) и отобразить их все в одном окне \ename{Line Char View}.
\end{itemize}
